\begin{figure*}[t]
\centering
\scriptsize
\setlength{\fboxsep}{1em}
\setlength{\fboxrule}{1pt}
\fcolorbox{black}{white}{
\parbox{0.955\textwidth}{
\begin{alltt}


You will be provided instructions for completing a task followed by a task to complete.

Instructions:

1. Understand the question: Read the question carefully and understand what it is asking. The questions in this dataset are based on elementary-level science topics and require multi-step reasoning, use of additional common and commonsense knowledge.

2. Analyze the options: Look at each of the multiple-choice options. The correct answer may not always be the most obvious one. Use your understanding of the question and the options to reason through the best answer.

3. Apply knowledge: The dataset includes a collection of common knowledge facts. Use these facts, along with your own knowledge, to answer the question. 

4. Be aware of the challenge: The questions are designed to be challenging for both retrieval-based and word co-occurrence algorithms. Don't rely solely on keyword matching or simple retrieval methods to answer the questions.

5. Check your answer: Before finalizing your answer, check it against the question and the options to make sure it makes sense and answers the question correctly.

Topics to know: The dataset covers a wide range of topics in science, including but not limited to: astronomy, biology, chemistry, physics, and earth science. It also requires a good understanding of common sense and the ability to apply this knowledge to novel situations. 

Steps to answer the questions: 

1. Read the question carefully and understand what it is asking.

2. Look at each of the multiple-choice options and analyze them in relation to the question.

3. Use your knowledge and the common knowledge facts provided in the dataset to reason through the best answer.
4. Check your answer against the question and the options to make sure it makes sense and answers the question correctly.

Question: A person wants to start saving money so that they can afford a nice vacation at the end of the year. After looking over their budget and expenses, they decide the best way to save money is to

A. make more phone calls

B. quit eating lunch out

C. buy less with monopoly money

D. have lunch with friends

Follow the instructions to generate an explanation that reasons towards the correct answer to the task above. End the explanation with the correct answer. [/INST] Explanation: 

To answer this question, we need to understand the situation and the given options. The person wants to save money for a vacation, and they have looked over their budget and expenses. They want to know the best way to save money.

Option A, making more phone calls, does not seem relevant to saving money. Making phone calls may even cost money, depending on the person's phone plan.

Option B, quitting eating lunch out, could be a good way to save money. Eating out can be expensive, and cooking meals at home can be cheaper. This option aligns with the person's goal of saving money.

Option C, buying less with monopoly money, does not make sense. Monopoly money is not real money and cannot be used to buy things in the real world. This option does not help the person achieve their goal of saving money.
Option D, having lunch with friends, may not be the best option for saving money. Depending on the cost of the meal and the number of friends, this option could end up costing more money than cooking meals at home.

Therefore, the correct answer is B, quitting eating lunch out. This option aligns with the person's goal of saving money and is a practical solution to reducing expenses.

Answer: B

Therefore, the correct multiple choice label (just the letter) to the task is below.

Answer: B
\end{alltt}
}}
\caption{A complete running example of OpenBookQA.}
\label{fig:commonsense_dataset=openbookqa_example}
\end{figure*}
